\documentclass{beamer}
\usepackage{amsfonts,amsmath,oldgerm}
\usepackage{graphicx} % Required for inserting images
\usepackage{caption}
\usepackage{subcaption}
\usepackage{amsthm}
\usepackage{mathtools}
\usepackage{float}

\usetheme{dmpisa} % Here the call to the beamer theme
\usepackage{amssymb}
\usepackage{physics}
\usepackage{algorithm}
\usepackage{algpseudocode}

\newcommand{\prefactor}[2]{A_{#1,#2}}

\title[Softmin Energy Minimization]{Global Optimization via Softmin Energy Minimization}
\author[AA, MR, SS, VC]{Andrea Agazzi, Marco Romito, \emph{Samuele Saviozzi}, Vittorio Carlei}
\date{June 2026}

\begin{document}

\frame{\titlepage}

% Section 1: Introduction
\section{Introduction}

\begin{frame}
\frametitle{The Global Optimization Problem}
\begin{itemize}
  \item \textbf{Goal}: Find global minimum of $\min_{x \in \mathbb{R}^d} f(x)$
  \item \textbf{Challenge}: Non-convex functions with multiple local minima
  \item Standard gradient descent gets trapped in local minima
  \item Need methods that balance exploration and exploitation
\end{itemize}
\pause
\vspace{0.5cm}
\textbf{Common approaches:}
\begin{itemize}
  \item \textbf{Particle Swarm Optimization}: multiple interacting particles, usually metaheuristic methods. (PSO, CBO, etc...)
  \item \textbf{Simulated Annealing}: gradient descent + noise (Langevin dynamics) that uses temperature annealing.
\end{itemize}
\end{frame}

\begin{frame}
\frametitle{Key Insight: Particle Cooperation}
Our approach combines:
\begin{enumerate}
  \item \textbf{Multiple particles}: parallel exploration
  \item \textbf{Gradient information}: exploit when available
  \item \textbf{Cooperation mechanism}: particles interact via Softmin Energy
\end{enumerate}
\pause

\vspace{0.5cm}
\textbf{Main innovation:} Follow the gradient of the \textbf{Softmin Energy}, a smooth approximation of the minimum value.

\end{frame}

% Section 2: Softmin Energy
\section{Softmin Energy Formulation}

\begin{frame}
\frametitle{The Softmin Energy Functional}

Given $n$ particles $x = (x_1, \ldots, x_n)$ with $x_i \in \mathbb{R}^d$, the Softmin Energy is:

\[
    J_\beta(x) = \sum_{i=1}^{n} \frac{\exp(-\beta f(x_i))}{\sum_{j=1}^{n} \exp(-\beta f(x_j))} f(x_i)
\]
\pause

\textbf{Interpretation:}
\begin{itemize}
  \item Weighted average of function values at particle positions
  \item Weight of particle $i$ depends on how low its function value is
  \item Parameter $\beta > 0$ controls smoothness:
  \begin{itemize}
    \item $\beta \to 0$: recovers average $\frac{1}{n}\sum_i f(x_i)$
    \item $\beta \to \infty$: approaches hard minimum $\min_i f(x_i)$
  \end{itemize}
\end{itemize}
\end{frame}

\begin{frame}
\frametitle{Gradient of Softmin Energy}
\begin{enumerate}
    \item Let $w_k = e^{-\beta f(x_k)}(\sum_j e^{-\beta f(x_j)})^{-1}$ denote the Gibbs weight of particle $x_k$
    \item Dynamic pre-factor $\prefactor{t}{k} = w_k\left[1 - \beta(f(x_k) - J_\beta(\mathbf{x}))\right]$.
\end{enumerate}
\pause
The gradient of $J_\beta(\mathbf{x})$ with respect to particle $x_k$ is given by:
\begin{equation}
    \nabla_{x_k} J_\beta(\mathbf{x}) =\prefactor{t}{k} \nabla f(x_k).
    \label{eq:softmin-gradient}
\end{equation}
The dynamic pre-factor $\prefactor{t}{k} = w_k\beta\left[1/\beta - f(x_k) + J_\beta(\mathbf{x})\right]$ dictates the effective direction of the gradient flow
\pause

\textbf{Key observation: }
\begin{enumerate}
    \item If $f(x_k) < J_\beta(\mathbf{x}) + 1/\beta$, $\prefactor{t}{k}$ is strictly positive, and the particle performs standard gradient descent on $f$.
    \item If $f(x_k) > J_\beta(\mathbf{x}) + 1/\beta$, the pre-factor becomes negative, reversing the flow.
\end{enumerate}
\end{frame}

\begin{frame}
\frametitle{Gradient Flow Dynamics}

\textbf{Deterministic dynamics:}
\begin{equation}
dx_t = -n \nabla J_{\beta_t}(x_t) dt
\end{equation}

\textbf{Stochastic dynamics (with exploration noise):}
\begin{equation}
dx_t = -n \nabla J_{\beta_t}(x_t) dt + \sqrt{2\sigma} dB_t
\end{equation}

where $B_t$ is Brownian motion and $\sigma > 0$ controls noise intensity. The parameter $\beta_t$ can decay over time.
\end{frame}

% Section 3: Theory
\section{Theoretical Results}

\subsection{Stability of Stationary Points}

\begin{frame}
\frametitle{Convergence for Strongly Convex Functions}

\textbf{Theorem 2.2:} For $\lambda$-strongly convex $f$ with min value 0:

At any stationary point $\bar{x}$:
\begin{itemize}
  \item At least $m \geq 1$ particles satisfy $\nabla f(x) = 0$ (at critical points)
  \item Remaining $n-m$ particles satisfy $f(x) = \frac{1}{\beta} + J_\beta(x)$
\end{itemize}
\pause

\textbf{Stability:}
\begin{itemize}
  \item If $m = n$: \textbf{stable}
  \item If $m = n-1$: Unstable saddle point $\Rightarrow$ exploration
\end{itemize}
\pause
\textbf{Implication:} The method finds the global minimum while maintaining exploratory behavior at suboptimal particles.
\end{frame}

% Section 4: Hitting Times
\subsection{Faster Barrier Crossing}

\begin{frame}{Exit Times from Local Minima: Intuition}
    To develop intuition, Consider a standard double well landscape
    \only<1>{
    \begin{figure}
        \centering
        \includegraphics[width=0.45\linewidth]{NIPS/Softmin/images/Screenshot 2026-01-20 125125.png}
    \end{figure}
    }
    \only<2>{
    \begin{figure}
        \centering
        \includegraphics[width=0.45\linewidth]{NIPS/Softmin/Slides/cat2.jpg}
    \end{figure}
    }
    \only<3>{
    \begin{figure}
        \centering
        \includegraphics[width=0.45\linewidth]{NIPS/Softmin/Slides/cat1.jpg}
    \end{figure}
    }
    \only<4>{
    \begin{figure}
        \centering
        \includegraphics[width=0.45\linewidth]{NIPS/Softmin/Slides/cat3.jpg}
    \end{figure}
    }
\end{frame}

\begin{frame}{Exit Times from Local Minima}
    By the Freidlin-Wentzell Large Deviation Principle (LDP), the exponential scale of the exit time is governed by the minimum action path required to reach the boundary $\partial D$. Let $V_1$ be the exact one-particle Freidlin-Wentzell quasipotential
    
    We define the explicit Softmin barrier bound $\overline{V}_\beta$ as:
    \begin{equation}
        \overline{V}_\beta := \max\left\{ \frac{n}{2(n-1)\beta}, \frac{\beta}{2}\left(f^\ast-\frac1\beta\right)_+^2 \right\}
    \end{equation}
\end{frame}

\begin{frame}{Exit Times from Local Minima: Result}
    \begin{block}{Theorem 3.3}
        Assume that $x^\ast$ is the minimal saddle point on the boundary $\partial D$. Then there exists $n_0\in\mathbb{N}$ such that, for all $n\ge n_0$
        \[
            \lim_{\sigma\downarrow0} \sigma\log \mathbb{E}_{\bar x}\tau_k^\sigma = V_1.
        \]
        Moreover, the probability that exactly one particle exits $D$ during the transition tends to $1$ as $\sigma\downarrow0$. Furthermore, the exponential scale is bounded by:
        \[
            V_1\le \overline{V}_\beta.
        \]
        Consequently, for every $\eta>0$, there exists $\sigma_0>0$ such that, for all $0<\sigma<\sigma_0$,
        \[
            \mathbb{E}_{\bar x}\tau_k^\sigma \le \exp\left(\frac{\overline{V}_\beta+\eta}{\sigma}\right).
        \]
    \end{block}
\end{frame}

\begin{frame}\frametitle{Sketch of the proof}
    \textbf{Establish channel separation:} we prove that multi-particle saddles ($V_q$ for $q \ge 2$) have strictly larger action than the one-particle saddle $V_1$.\\
    \pause
    \textbf{Compute the upper bound $\overline{V}_\beta$:} we evaluate a specific channel where $n-1$ particles are frozen at $x_{\min}$ while a single active particle escapes. This provides a strict upper bound. \\
    \pause
    This 1D channel structurally decomposes into two bottlenecks:
    \begin{itemize}
        \item \textbf{Ejection:} Bounding the active particle's Gibbs weight yields a linear majorant for the effective learning rate. Integrating this majorant caps the ejection action at $n/(2(n-1)\beta)$.
        \pause
        \item \textbf{Descent:} Integrating the maximal opposing gradient provides a strict bound on the descent action of $(\beta/2)(f^* - 1/\beta)_+^2$.
    \end{itemize}
    Taking the maximum of these two sequential bottlenecks yields $V_1 \le \overline{V}_\beta$.
\end{frame}

\subsection{Speed of Convergence}
\begin{frame}

\frametitle{Convergence to Local Minima}

\textbf{Theorem 3.7:} For locally strongly convex $f$, if particle starts in "strongly minimizing regime":
$$\|x_k(t) - x^*\| \leq \|x_k(0) - x^*\| e^{-\lambda t}$$

Convergence rate: $O(e^{-\lambda t})$ where $\lambda$ is strong convexity constant.
\pause

\vspace{0.5cm}
\textbf{Comparison with Langevin dynamics:}
\begin{itemize}
  \item Both achieve exponential convergence
  \item Softmin method has prefactor $\prefactor{t}{k} > 1$ vs. standard prefactor $1$
\end{itemize}
\end{frame}

% Section 5: Numerical Experiments
\section{Numerical Experiments}

\begin{frame}
\frametitle{Benchmark 1: 1D Double Well}

\begin{figure}[htbp]
    \centering
    % Left Image
    \begin{subfigure}[b]{0.45\textwidth}
        \centering
        \includegraphics[width=\textwidth]{NIPS/Softmin/images/1d_double_well_ls.png}
    \end{subfigure}
    \hfill % This pushes the two images apart
    % Right Image
    \begin{subfigure}[b]{0.48\textwidth}
        \centering
        \includegraphics[width=\textwidth]{NIPS/Softmin/images/1dim_decay_big.png}
    \end{subfigure}
    \label{fig:1d_double_well}
\end{figure}

\textbf{Result:} Softmin method (fixed $\beta$ and decaying $\beta$) escapes local minima faster for high barriers.
\end{frame}

\begin{frame}
\frametitle{Benchmark 2: 2D Quadruple Well}

\begin{figure}[htbp]
    \centering
    % Left Image
    \begin{subfigure}[b]{0.38\textwidth}
        \centering
        \includegraphics[width=\textwidth]{NIPS/Softmin/images/2d_double_well_ls.png}
    \end{subfigure}
    \hfill % This pushes the two images apart
    % Right Image
    \begin{subfigure}[b]{0.48\textwidth}
        \centering
        \includegraphics[width=\textwidth]{NIPS/Softmin/images/2dim_big.png}
    \end{subfigure}
    \label{fig:2d_double_well}
\end{figure}

\textbf{Result:} Superior performance at high barriers. Transition times show different scaling law compared to Simulated Annealing.
\end{frame}

\begin{frame}
\frametitle{Benchmark 3: Ackley Function}

\begin{figure}[htbp]
    \centering
    % Left Image
    \begin{subfigure}[b]{0.35\textwidth}
        \centering
        \includegraphics[width=\textwidth]{NIPS/Softmin/images/Ackley_function_ls.png}
    \end{subfigure}
    \hfill % This pushes the two images apart
    % Right Image
    \begin{subfigure}[b]{0.48\textwidth}
        \centering
        \includegraphics[width=\textwidth]{NIPS/Softmin/images/high_dim.png}
    \end{subfigure}
    \label{fig:ackley}
\end{figure}

\textbf{Result:} Softmin method scales better with the dimension.
\end{frame}

% Section 6: Summary
\section{Summary}

\begin{frame}
\frametitle{Key Contributions}

\begin{enumerate}
  \item \textbf{Novel method:} Gradient-based particle swarm via Softmin Energy
  
  \item \textbf{Theoretical guarantees:}
  \begin{itemize}
    \item Convergence to global minimum for strongly convex functions
    \item Exponential speedup in barrier crossing (lower effective barriers)
    \item Exponential convergence rates to local minima
  \end{itemize}
  
  \item \textbf{Practical advantages:}
  \begin{itemize}
    \item Particle cooperation speeds up potential barriers crossing
    \item Outperforms Particle Swarm methods in various regimes on benchmark problems
  \end{itemize}
\end{enumerate}
\end{frame}

\begin{frame}
\frametitle{Why This Works}

\textbf{The mechanism:}
\begin{itemize}
  \item Softmin Energy creates cooperative interaction between particles
  \item Good particles (low function values) converge to minima faster
  \item Bad particles are pushed to explore and escape local minima
\end{itemize}

\vspace{0.5cm}
\textbf{Temperature annealing:}
\begin{itemize}
  \item Parameter $\beta$ plays role analogous to temperature
  \item Start with $\beta$ big for broad exploration and decrease over time
\end{itemize}

\end{frame}

\begin{frame}
\frametitle{Future Directions}

\begin{itemize}
  \item Extend proofs to weakly convex and non-convex functions
  \item Adaptive parameter tuning for $\beta$
  \item Application to machine learning (training neural networks)
  \item Mean-field analysis of the infinite particle limit
\end{itemize}

\end{frame}

\begin{frame}
\frametitle{Thank You For The Attention!}

\begin{figure}
    \centering    \includegraphics[width=0.25\linewidth]{NIPS/Softmin/images/Qq1hbauAiniaN4V.png}
\end{figure}

\begin{itemize}
  \item Paper: \textit{Global Optimization via Softmin Energy Minimization}
  \item Contact: samuele.saviozzi@phd.unipi.it
\end{itemize}
\end{frame}

\end{document}
